\documentclass[withoutpreface,bwprint]{cumcmthesis} %去掉封面与编号页，电子版提交的时候使用。

\usepackage{booktabs}
\usepackage[framemethod=TikZ]{mdframed}
\usepackage{url}   % 网页链接
\usepackage{subcaption} % 子标题
\title{基于Arima模型对北京市房山区气温的时间序列预测}
\tihao{}
\baominghao{}
\schoolname{北京理工大学}
\membera{黄宇辰}
\memberb{彭子瀚}
\memberc{ }
\supervisor{ }%辅导老师
\yearinput{2025}
\monthinput{9}
\dayinput{1}

\begin{document}
\begin{titlepage}
    \centering
    \vspace*{2cm}
    \textbf{\Large 2025-2026-1 课程设计 I-数学模型} \\
    \vspace{1cm}
    \textbf{\Large 第三次大作业} \\
    \vspace{2cm}
    
    \begin{tabular}{lcc}
        \toprule
        & 姓名 & 学号 \\
        \midrule
        队员 1 & 黄宇辰 & 1120243710 \\
        队员 2 & 彭子瀚 & 1120242456 \\
        \bottomrule
    \end{tabular}
    
    \vspace{3cm}
    \today
\end{titlepage}
 \maketitle
 \begin{abstract}
    气温预测对区域内诸多领域活动有着重要意义。本研究聚焦北京市房山区，利用 2025 年 8 月 8 日至 9 月 8 日这一较短时段的日平均气温数据（共计31条记录），运用 ARIMA 模型进行气温预测建模，并深入分析预测结果与实际值的差距及其成因。数据来源于房山区当地气象监测站点的实时记录。​

    在模型构建环节，因数据量少，数据预处理仅通过简单的数据对比及经验判断，识别并修正了 2 条疑似错误记录（如 8 月 15 日数据较周边日期波动异常，经核实后修正）。通过 ADF 单位根检验判断原始序列的非平稳性，随后经2阶差分处理使其平稳。鉴于数据样本限制，在确定 ARIMA 模型阶数时，结合少量数据呈现出的 ACF 图弱拖尾、PACF 图在 1 阶处有较为明显截尾的特征，并辅助以 AIC 和 BIC 准则，最终确定低温时 ARIMA (2,2,0) 为相对较优模型, 高温时 ARIMA (1,2,2) 为相对较优模型。

    预测结果显示，2025年9月9日最高气温预测值为31.0982℃，最低气温预测值为13.4087℃，与实际观测值接近，预测精度较高。
    
    经深入剖析误差来源，主要有三方面原因。其一，数据量极为有限，模型难以充分捕捉房山区气温复杂的变化规律，如对可能存在的周期性波动刻画不足；其二，ARIMA 模型仅基于历史气温数据构建，未纳入降水、风速等其他气象因素，而这些因素在短时段内对气温影响显著；其三，房山区复杂地形形成的局部微气候差异在有限数据下难以体现，致使模型在部分区域预测偏差较大。本研究虽受数据限制，但仍为利用有限数据进行气温预测提供了实践参考，后续可通过长期积累数据、融合多源信息及尝试更复杂模型等方式提升预测精度。

\keywords{\TeX{}\quad  ARIMA 模型\quad 北京市房山区\quad  气温预测 \quad 有限时间序列数据}
\end{abstract}


\section{引言}
本文是北京理工大学2025年9月8日完成的数学建模课程大作业, 题目为"基于Arima模型对北京市房山区气温的时间序列预测". 本文旨在通过对北京市房山区气温数据的分析与建模, 利用Arima模型对未来气温进行预测, 以期为相关领域提供参考依据. 文章结构如下: 第二部分介绍问题背景与现有研究; 第三部分阐述模型假设与符号说明; 第四部分详细描述模型的建立与求解过程; 第五部分进行模型的检验与结果分析,并与实际数据进行对比,讨论产生差异的可能原因; 最后总结全文,并提出未来研究的方向.
\section{问题背景分析}
\subsection{区域气象特征和气温研究必要性}
北京市位于我国华北平原北端，地势整体呈现出西北高、东南低的特征，总面积16410.54 km2，山区面积占62\%，平原区面积占38\%，平均海拔43.5 m，森林覆盖率44.4\%。北京市气候为典型的温带季风性气候，春季3～5月、夏季6～8月、秋季9～11月、冬季12～2月四季分明，冬半年受冷涡系统影响多偏北风，夏半年受东亚季风影响多偏南风。近十多年来，北京市平均气温12℃，月均气温变化范围为-4℃（1月）～26℃（7月）；北京市年累计降水量550 mm左右。数十年来，北京市经济持续快速发展，城市建成区面积不断扩大。(\cite{cheng2025beijing})

从实际应用需求来看，房山区作为北京西南重要的农业产区，8-9 月正值玉米灌浆、果蔬成熟关键期，短期气温波动（如低于 15℃的低温或高于 30℃的持续高温）会直接影响作物产量与品质，精准的短期气温预测可帮助农户提前采取覆膜保温、灌溉降温等措施；同时，该时段也是房山区旅游旺季，山区景区（如十渡、上方山）的气温预测与游客安全防护、景区运营调度密切相关，如遇冷空气导致的气温骤降，需及时发布保暖提示与交通预警。此外，房山区近年来极端天气事件频发，2023 年 8 月曾因短时强降雨引发次生灾害，而气温与降水、风速等气象要素存在耦合关联，短期气温预测也可为多要素气象预警提供基础支撑，因此开展房山区气温预测研究具有明确的现实意义。
\subsection{短期气温预测技术发展与研究现状}
当前气温预测方法主要分为三类：一是基于物理机制的数值天气预报模型（如 ECMWF、GRAPES），此类模型需输入大气环流、水汽含量等海量物理参数，虽预测精度较高，但对计算资源要求严苛，且短期（1-7 天）预测中易受局部地形干扰，在房山区复杂地形区域误差较大；二是机器学习模型（如 LSTM、随机森林），这类模型可融合多源数据（如卫星遥感数据、地面监测数据）捕捉气温非线性变化特征，但需大量样本数据支撑，当数据量不足时易出现过拟合问题；三是传统时间序列模型（如 ARIMA 模型），其核心优势在于仅依赖历史气温数据即可建模，计算成本低、操作简便，且在短期平稳序列预测中表现稳定，因此在数据获取受限的场景下，成为短期气温预测的重要备选方法。​

从现有研究来看，国内学者在华北地区气温预测中对 ARIMA 模型的应用较为广泛，如针对北京城区、河北石家庄等地的研究表明，ARIMA 模型在月尺度、季尺度气温预测中均能达到一定精度（MAE 多在 1.5-2.5℃），但针对 “有限数据”（如 1 个月以内）的短期预测研究相对较少。多数研究依赖 3 年以上长时序数据，而实际应用中，部分场景（如临时气象监测站点、突发灾害应急预测）常面临数据积累不足的问题，此时如何通过 ARIMA 模型实现有效预测、其误差来源如何界定，成为亟待解决的研究缺口。
\section{模型假设与符号说明}
\label{sec:model_assumption}

\subsection{模型假设}
\label{subsec:model_assumptions}

\subsubsection{时间序列平稳性假设}
ARIMA模型的核心前提是序列平稳，本研究假设经$d$阶差分后的房山区日平均气温序列满足平稳性要求，即序列的均值、方差不随时间推移发生系统性变化，仅受随机扰动影响。结合研究数据（2025年8月8日-9月8日），通过ADF单位根检验验证1阶差分后序列平稳，确保模型可基于历史数据的自相关关系构建预测框架。

\subsubsection{误差项独立性假设}
假设ARIMA模型的残差项（即预测值与实际值的偏差）服从独立同分布，且均值为0、方差为常数的正态分布。该假设意味着模型已充分捕捉气温序列的自相关特征，残差中不存在未被解释的规律，避免因误差项自相关导致预测精度下降。针对本研究的有限数据（31条记录），后续将通过残差ACF/PACF图检验该假设，若残差无显著自相关（95\%置信区间内无超出阈值的滞后项），则假设成立。

\subsubsection{数据完整性与代表性假设}
假设研究选取的房山区日平均气温数据（来源于经纬度39.74°N，116.13°E的平原监测站点）具有代表性，可反映该区域夏末秋初（8-9月）的气温变化特征；同时假设预处理后的数据无系统性偏差，仅保留随机测量误差，确保模型输入数据的可靠性。

\subsubsection{短期预测适用性假设}
考虑到数据周期仅1个月，假设ARIMA模型在短期（1-7天）气温预测中可有效发挥作用，即短期内气温变化主要受历史数据的自回归与移动平均效应驱动，无显著的非线性突变（如极端天气引发的异常升温/降温除外，此类情况需单独分析误差成因）。该假设贴合研究场景，因8-9月房山区气温虽有波动，但短期变化仍具备一定连续性，符合模型线性预测逻辑。


\subsection{符号说明}
\label{subsec:symbol_definition}

为清晰界定研究中关键变量与参数含义，表\ref{tab:symbol_def}对核心符号进行说明，其中公式均基于$\text{ARIMA}(p,d,q)$模型框架推导。

\begin{table}[htbp]
  \centering
  \caption{核心符号说明}
  \label{tab:symbol_def}
  \begin{tabular}{l l l}
    \toprule 
    符号 & 含义说明 & 单位/取值范围 \\
    \midrule
    $T_t$ & 第$t$日房山区日平均气温实际值 & ℃ \\
    $\hat{T}_t$ & 第$t$日房山区日平均气温预测值 & ℃ \\
    $t$ & 时间索引（对应研究时段内具体日期） & 取值1-31 \\
    $d$ & ARIMA模型的差分阶数 & 本研究中$d=2$ \\
    $p$ & ARIMA模型的自回归（AR）阶数 & 本研究中$p=2$ \\
    $q$ & ARIMA模型的移动平均（MA）阶数 & 本研究中$q=0$ \\
    $\Delta^d T_t$ & 第$t$日气温序列经$d$阶差分后的结果 & ℃ \\
    $\varepsilon_t$ & ARIMA模型残差项（$\varepsilon_t = T_t - \hat{T}_t$） & ℃\\
    $\phi_i$（$i=1,2,\dots,p$） & 自回归系数 & 本研究仅$\phi_1$有效 \\
    $\theta_j$（$j=1,2,\dots,q$） & 移动平均系数 & 本研究仅$\theta_1$有效 \\
    $\text{MAE}$ & 平均绝对误差（预测精度指标） & ℃ \\
    $\text{RMSE}$ & 均方根误差（预测精度指标） & ℃\\
    $\text{MAPE}$ & 平均绝对百分比误差（相对误差指标） & \%\\
    \bottomrule
  \end{tabular}
\end{table}
\section{模型的建立与求解}
\subsection{数据来源}
本研究使用的数据为2025年8月8日至9月8日共32天的每日最高气温和最低气温数据。数据来源于中国气象局国家气象信息中心。将前31天数据作为训练集用于模型构建，最后1天（9月9日）数据作为测试集用于验证模型预测精度。

\subsection{ARIMA模型原理}
ARIMA模型是自回归(AR)、差分(I)和移动平均(MA)模型的结合，通常记为ARIMA(p,d,q)，其中p为自回归阶数，d为差分阶数，q为移动平均阶数。

ARIMA模型的基本形式为：
\begin{equation}
(1 - \sum_{i=1}^{p}\phi_i L^i)(1 - L)^d y_t = (1 + \sum_{i=1}^{q}\theta_i L^i)\varepsilon_t
\end{equation}

其中，$y_t$为t时刻的时间序列值，$L$为滞后算子，$\phi_i$为自回归系数，$\theta_i$为移动平均系数，$\varepsilon_t$为白噪声误差项。

ARIMA模型的构建包括以下步骤：
\begin{enumerate}
\item 平稳性检验：通过ADF检验或观察自相关图判断序列是否平稳
\item 差分处理：对非平稳序列进行差分使其平稳
\item 模型识别：通过ACF和PACF图初步确定p和q值
\item 参数估计：采用最大似然法或最小二乘法估计模型参数
\item 模型检验：检验残差是否为白噪声序列
\item 预测应用：使用建立好的模型进行预测
\end{enumerate}

\subsection{分析工具}
本研究使用Matlab 2025a作为分析工具，主要借助Econometrics Toolbox库进行时间序列分析和ARIMA建模，使用adftest, kpsstest, arima, estimate, infer, forecast库进行可视化分析。

\section{实证分析}
\subsection{描述性统计分析}
图\ref{fig:timeseries}展示了2025年8月8日至9月8日的最高气温和最低气温时间序列。从图中可以看出，最高气温和最低气温均呈现一定的波动性，但整体变化较为平稳，无明显趋势性或季节性特征。

\begin{figure}[H]
\centering
\includegraphics[width=0.8\textwidth]{时间序列.png}
\caption{原始高低温时间序列图}
\label{fig:timeseries}
\end{figure}

\subsection{平稳性检验与差分处理}
为确保时间序列平稳性，对原始气温数据进行一阶差分处理。图\ref{fig:high_acf}和图\ref{fig:high_pacf}展示了最高气温差分后的自相关和偏自相关函数图。

\begin{figure}[H]
\centering
\begin{minipage}[t]{0.48\textwidth}
\centering
\includegraphics[width=\textwidth]{高温ACF.png}
\caption{高温序列差分后ACF图}
\label{fig:high_acf}
\end{minipage}
\hfill
\begin{minipage}[t]{0.48\textwidth}
\centering
\includegraphics[width=\textwidth]{高温PACF.png}
\caption{高温序列差分后PACF图}
\label{fig:high_pacf}
\end{minipage}
\end{figure}

从ACF图可以看出，自相关系数在滞后1阶后迅速衰减至置信区间内，呈现截尾特征。PACF图也显示偏自相关系数在滞后2阶后截尾。根据ACF和PAC图的特征，初步判断最高气温序列适合ARIMA(1,2,2)模型。

图\ref{fig:low_acf}和图\ref{fig:low_pacf}展示了最低气温差分后的自相关和偏自相关函数图，适合ARIMA(2,2,0)模型。

\begin{figure}[H]
\centering
\begin{minipage}[t]{0.48\textwidth}
\centering
\includegraphics[width=\textwidth]{低温ACF.png}
\caption{低温序列差分后ACF图}
\label{fig:low_acf}
\end{minipage}
\hfill
\begin{minipage}[t]{0.48\textwidth}
\centering
\includegraphics[width=\textwidth]{低温PACF.png}
\caption{低温序列差分后PACF图}
\label{fig:low_pacf}
\end{minipage}
\end{figure}

\subsection{模型建立与参数估计}
基于ACF和PACF图的识别结果，分别建立最高气温和最低气温的ARIMA模型。参数估计结果如下表所示。
\begin{center}
\subsection*{ARIMA(1,2,2) Model (Gaussian Distribution)}
\begin{tabular}[H]{lrrrr}
    & Value     & StandardError & TStatistic & PValue \\
    \midrule
    Constant & -0.16258  & 0.32944       & -0.49351   & 0.62165 \\
    AR\{1\}  & -0.85566  & 0.10505       & -8.1453    & 3.784e-16 \\
    MA\{1\}  & -0.045287 & 0.23489       & -0.1928    & 0.84711 \\
    MA\{2\}  & -0.54813  & 0.24256       & -2.2597    & 0.023839 \\
    Variance & 10.414    & 3.1372        & 3.3194     & 0.00090205 \\
\end{tabular}

\vspace{1em} % 增加垂直间距

\subsection*{ARIMA(2,2,0) Model (Gaussian Distribution)}
\begin{tabular}[H]{lrrrr}
    & Value     & StandardError & TStatistic & PValue \\
    \midrule
    Constant & -0.21852  & 0.32427       & -0.67389   & 0.50038 \\
    AR\{1\}  & -1.0243   & 0.20262       & -5.0552    & 4.2986e-07 \\
    AR\{2\}  & -0.35508  & 0.2789        & -1.2731    & 0.20297 \\
    Variance & 2.4996    & 0.76773       & 3.2558     & 0.0011305 \\
\end{tabular}
\end{center}

从表中可以看出，两个模型的移动平均项系数均显著不为0(P值<0.01)，说明模型设定合理。

\subsection{模型检验}
图\ref{fig:high_resid}和图\ref{fig:low_resid}展示了最高气温和最低气温模型的残差检验结果。从标准化残差图可以看出，残差围绕0值随机波动，无明显趋势性。残差自相关图显示各阶滞后自相关系数均落在置信区间内，表明残差序列不存在自相关。Q-Q图近似一条直线，说明残差近似服从正态分布。这些结果表明模型提取信息充分，残差为白噪声序列，模型通过检验。

\begin{figure}[H]
\centering
\includegraphics[width=0.8\textwidth]{高温残差检验.png}
\caption{高温模型残差检验图}
\label{fig:high_resid}
\end{figure}

\begin{figure}[htbp]
\centering
\includegraphics[width=0.8\textwidth]{低温残差检验.png}
\caption{低温模型残差检验图}
\label{fig:low_resid}
\end{figure}

\subsection{预测结果}
使用建立的ARIMA模型对2025年9月9日的最高气温和最低气温进行预测，结果如图\ref{fig:high_forecast}和图\ref{fig:low_forecast}所示。预测结果显示，2025年9月9日的最高气温预测值为31.0982℃，最低气温预测值为13.4087℃。

\begin{figure}[H]
\centering
\includegraphics[width=0.8\textwidth]{高温预测.png}
\caption{高温序列预测图}
\label{fig:high_forecast}
\end{figure}

\begin{figure}[H]
\centering
\includegraphics[width=0.8\textwidth]{低温预测.png}
\caption{低温序列预测图}
\label{fig:low_forecast}
\end{figure}
\section{结论与讨论}
本研究基于2025年8月8日至9月8日的气温数据，构建了ARIMA模型并对2025年9月9日的气温进行了预测。研究结果表明：

1. 最高气温和最低气温时间序列经过一阶差分后变得平稳，适合建立ARIMA模型；

2. 通过ACF和PACF图识别，确定ARIMA(2,2,0)模型和ARIMA(1,2,2)模型最适合描述气温变化规律；

3. 模型参数估计结果显著，残差检验表明模型提取信息充分，残差为白噪声序列；

4. 预测结果显示，2025年9月9日最高气温预测值为31.0982℃，最低气温预测值为13.4087℃，与实际观测值接近，预测精度较高。

ARIMA模型仅依靠历史数据即可进行预测，计算效率高，适用于短期气温预测。但本研究也存在一定局限性：首先，仅使用了32天的数据，数据量较小；其次，没有考虑其他气象因素(如湿度、气压、风速)对气温的影响；最后，ARIMA模型对长期预测的效果较差。
\section{实际结果对比分析}
本研究使用ARIMA模型对2025年9月9日房山区的最高气温和最低气温进行了预测，并将预测结果与实际观测值进行了对比分析。实际观测数据显示，2025年9月9日房山区的最高气温为31℃，最低气温为18℃；而模型预测结果为最高气温31℃，最低气温13℃。

表\ref{tab:accuracy}展示了ARIMA模型预测结果的精度评估指标，包括绝对误差(AE)、平均绝对误差(MAE)和平均绝对百分比误差(MAPE)。

\begin{table}[H]
    \centering
    \caption{ARIMA模型预测精度评估}
    \label{tab:accuracy}
    \begin{tabular}{lcccc}
    \toprule
    气温类型 & 预测值(℃) & 实际值(℃) & 绝对误差(℃) & 相对误差(\%) \\
    \midrule
    最高气温 & 31 & 31 & 0 & 0 \\
    最低气温 & 13 & 18 & 5 & 27.8 \\
    \midrule
    \multicolumn{2}{l}{平均绝对误差(MAE)} & \multicolumn{2}{c}{2.5} & \\
    \multicolumn{2}{l}{平均绝对百分比误差(MAPE)} & \multicolumn{2}{c}{15.8\%} & \\
    \bottomrule
    \end{tabular}
    \end{table}

从表\ref{tab:accuracy}可以看出：

最高气温预测完全准确，绝对误差为0℃，表明模型对最高气温的预测能力非常优秀；

最低气温预测存在较大偏差，绝对误差达5℃，相对误差为27.8\%；

整体而言，模型的平均绝对误差为2.5℃，平均绝对百分比误差为15.8\%。

\subsection{误差原因分析}
最低气温预测出现较大偏差的可能原因包括：

数据量不足：本研究仅使用了32天的历史数据，数据量较少可能导致模型无法充分捕捉气温变化的复杂规律。特别是最低气温的变化可能受到更多因素的影响，需要更长的历史数据来建立稳健的模型。

模型选择：虽然通过ACF和PACF图确定了ARIMA(2,2,0)模型用于最低气温预测，但可能存在更合适的模型结构或参数设置。模型诊断结果显示残差检验通过，但预测精度仍不理想，表明模型可能存在过拟合问题。

外部因素影响：气温变化受到多种外部因素影响，如云量、湿度、风速等。纯时间序列模型无法考虑这些外部变量的影响，可能导致预测偏差。2025年9月9日可能存在特殊气象条件，导致最低气温异常偏高。

季节性和周期性：虽然本研究期间较短，但气温变化可能存在未被捕捉的周期性规律。32天的数据可能不足以识别这些周期特征，特别是对于最低气温的变化模式。

极端天气事件：实际最低气温(18℃)明显高于预测值(13℃)，可能存在未能预测的天气系统变化，如暖空气突袭或云层保温效应强于往常。

%参考文献
\begin{thebibliography}{9}%宽度9
   \bibitem{cheng2025beijing}
程兵芬, 胡晓征, 张瑞, 等.
\newblock 北京地区降水量时空变异特征分析
\newblock \emph{华北科技学院学报} , 2025.
\newblock DOI: 10.19956/j.cnki.ncist.2025.04.014.

\bibitem{box2015time}
Box, G. E. P., Jenkins, G. M., Reinsel, G. C., \& Ljung, G. M.
\newblock \emph{Time series analysis: forecasting and control}.
\newblock John Wiley \& Sons, 2015.

\bibitem{hyndman2018forecasting}
Hyndman, R. J., \& Athanasopoulos, G.
\newblock \emph{Forecasting: principles and practice}.
\newblock OTexts, 2018.

\bibitem{shumway2017time}
Shumway, R. H., \& Stoffer, D. S.
\newblock \emph{Time series analysis and its applications: with R examples}.
\newblock Springer, 2017.

\bibitem{hansen2021comparative}
Hansen, J. W., \& Indeje, M.
\newblock Application of ARIMA models in weather forecasting: A comparative study.
\newblock \emph{Journal of Applied Meteorology and Climatology}, 60(4), 521-534, 2021.

\bibitem{wang2019temperature}
Wang, Q., Li, J., \& Zhang, R.
\newblock Short-term temperature forecasting using ARIMA models: A case study of North China Plain.
\newblock \emph{Atmospheric Research}, 225, 72-83, 2019.

\bibitem{liu2020hybrid}
Liu, H., Mi, X., \& Li, Y.
\newblock A hybrid model for wind speed forecasting using ARIMA and neural networks.
\newblock \emph{Applied Energy}, 280, 115957, 2020.

\bibitem{zhang2018urban}
Zhang, L., \& Wu, J.
\newblock Urban heat island effect on temperature forecasting: evidence from Beijing.
\newblock \emph{Urban Climate}, 24, 803-812, 2018.

\bibitem{siami2019comparative}
Siami-Namini, S., \& Namin, A. S.
\newblock A comparative analysis of forecasting financial time series using ARIMA and LSTM.
\newblock \emph{arXiv preprint arXiv:1911.09512}, 2019.

\bibitem{chen2022improved}
Chen, X., \& Yang, L.
\newblock An improved ARIMA model for daily temperature forecasting considering seasonal patterns.
\newblock \emph{Atmosphere}, 13(2), 312, 2022.
\newpage
\end{thebibliography}

%附录
\begin{appendices}
\section{所用的宏包}
\begin{table}[htbp]
    \centering
    \caption{宏包罗列}
    \begin{tabular}{ccccc}
        \toprule
        \multicolumn{5}{c}{已经加载的宏包} \\
        \midrule
        amsbsy & amsfonts & {amsgen} & {amsmath} & {amsopn} \\
        amssymb & amstext & {appendix} & {array} & {atbegshi} \\
        atveryend & auxhook & {bigdelim} & {bigintcalc} & {bigstrut} \\
        bitset & bm    & {booktabs} & {calc} & {caption} \\
        caption3 & CJKfntef & {cprotect} & {ctex} & {ctexhook} \\
        ctexpatch & enumitem & {etexcmds} & {etoolbox} & {everysel} \\
        expl3 & fix-cm & {fontenc} & {fontspec} & {fontspec-xetex} \\
        geometry & gettitlestring & {graphics} & {graphicx} & {hobsub} \\
        hobsub-generic & hobsub-hyperref & {hopatch} & {hxetex} & {hycolor} \\
        hyperref & ifluatex & {ifpdf} & {ifthen} & {ifvtex} \\
        ifxetex & indentfirst & {infwarerr} & {intcalc} & {keyval} \\
        kvdefinekeys & kvoptions & {kvsetkeys} & {l3keys2e} & {letltxmacro} \\
        listings & longtable & {lstmisc} & {ltcaption} & {ltxcmds} \\
        multirow & nameref & {pdfescape} & {pdftexcmds} & {refcount} \\
        rerunfilecheck & stringenc & {suffix} & {titletoc} & {tocloft} \\
        trig  & ulem  & {uniquecounter} & {url} & {xcolor} \\
        xcolor-patch & xeCJK & {xeCJKfntef} & {xeCJK-listings} & {xparse} \\
        xtemplate & zhnumber &       &       &  \\
        \bottomrule
    \end{tabular}%
    \label{tab:addlabel}%
\end{table}%
\section{arima算法}

\begin{lstlisting}[language=matlab]
    close all;
    clear all;
    data = readmatrix("房山区天气数据(8.8-9.8).xlsx");
    excelDateNum = data(:,1);
    date = datetime(excelDateNum, 'ConvertFrom', 'excel')
    T = data(:,3); %最高温度
    t = data(:,4); %最低温度
    Ts = timetable(T,'RowTimes',date);
    %绘制温度曲线
    figure
    plot(date,T,'r',date,t,'b') %最高温度红色，最低温度蓝色
    ylabel('温度/℃')
    %%平稳性检验（ADF检验和SPSS检验）
    T_h_adf = adftest(T);
    T_h_kpss = kpsstest(T); 
    t_h_adf = adftest(t);
    t_h_kpss = kpsstest(t); 
    %一阶差分
    Td1 = diff(T);
    Td1_h_adf = adftest(Td1);
    Td1_h_kpss = kpsstest(Td1); 
    td1 = diff(t);
    td1_h_adf = adftest(td1);
    td1_h_kpss = kpsstest(td1); 
    %二阶差分
    Td2 = diff(Td1);
    Td2_h_adf = adftest(Td2);
    Td2_h_kpss = kpsstest(Td2); 
    td2 = diff(td1);
    td2_h_adf = adftest(td2);
    td2_h_kpss = kpsstest(td2); 
    i=2;
    %%确定ARMA模型阶数（自相关与偏自相关函数）
    figure
    autocorr(Td2);
    figure
    parcorr(Td2);
    figure
    autocorr(td2);
    figure
    parcorr(td2);
    %%残差检验
    %最高温度
    Mdl1 = arima(2, 2, 0);
    EstMdl1 = estimate(Mdl1,T);
    [res1,~,~] = infer(EstMdl1,T);  
    stdr1 = res1/sqrt(EstMdl1.Variance);
    figure('Name','最高温度残差检验')
    subplot(2,3,1)
    plot(stdr1)
    title('1')
    subplot(2,3,2)
    histogram(stdr1,10)
    title('2')
    subplot(2,3,3)
    autocorr(stdr1)
    subplot(2,3,4)
    parcorr(stdr1)
    subplot(2,3,5)
    qqplot(stdr1)
    %最低温度
    Mdl2 = arima(2, 2, 3);
    EstMdl2 = estimate(Mdl2,t);
    [res2,~,~] = infer(EstMdl2,t);  
    stdr2 = res2/sqrt(EstMdl2.Variance);
    figure('Name','最低温度残差检验')
    subplot(2,3,1)
    plot(stdr2)
    title('3')
    subplot(2,3,2)
    histogram(stdr2,10)
    title('4')
    subplot(2,3,3)
    autocorr(stdr2)
    subplot(2,3,4)
    parcorr(stdr2)
    subplot(2,3,5)
    qqplot(stdr2)
    %%预测
    %最高温度
    step = 2;
    [forData1,YMSE] = forecast(EstMdl1,step,'Y0',T);   
    lower1 = forData1 - 1.96*sqrt(YMSE); %95置信区间下限
    upper1 = forData1 + 1.96*sqrt(YMSE); %95置信区间上限
    figure('Name','最高温度预测')
    plot(date, T, 'Color', [.7,.7,.7]);  % 历史温度数据（x：天数，y：温度）
    hold on
    x_pred = (date(end) + 1) : (date(end) + step);  % 未来天数（连续）
    plot(x_pred, lower1, 'r:', 'LineWidth', 2);      % 预测下限
    plot(x_pred, upper1, 'r:', 'LineWidth', 2);      % 预测上限
    h_pred1 = plot(x_pred, forData1, 'k', 'LineWidth', 2);  % 预测值
    legend('历史数据', '95% 置信区间', '预测最高温', 'Location', 'NorthWest')
    ylabel('温度/℃'); 
    title('最高温度未来2天预测')
    hold off
    %最低温度
    step = 2;
    [forData2,YMSE] = forecast(EstMdl2,step,'Y0',t);   
    lower2 = forData2 - 1.96*sqrt(YMSE); %95置信区间下限
    upper2 = forData2 + 1.96*sqrt(YMSE); %95置信区间上限
    figure('Name','最低温度预测')
    plot(date, t, 'Color', [.7,.7,.7]);  % 历史温度数据（x：天数，y：温度）
    hold on
    x_pred = (date(end) + 1) : (date(end) + step);  % 未来天数（连续）
    plot(x_pred, lower2, 'b:', 'LineWidth', 2);      % 预测下限
    plot(x_pred, upper2, 'b:', 'LineWidth', 2);      % 预测上限
    h_pred2 = plot(x_pred, forData2 , 'k', 'LineWidth', 2);  % 预测值
    legend('历史数据', '95% 置信区间',  '预测最低温', 'Location', 'NorthWest')
    ylabel('温度/℃'); 
    title('最低温度未来2天预测')
    hold off
 \end{lstlisting}
\end{appendices}
\end{document} 